\section{Компилятор Go}
\label{sec:Chapter4} \index{Chapter4}

\subsection{Существующие реализации компилятора Go и проблема векторизации}

Основным (и официально поддерживаемым) компилятором языка Go является \texttt{cmd/compile}~\cite{golang_compile_doc, golang_compile_readme}
 --- самодостаточный компилятор, ориентированный на высокую скорость сборки. 
 Именно он является основным в экосистеме Go. 
 Однако существует и несколько альтернативных реализаций,
 среди которых наиболее заметна \texttt{gollvm}~\cite{gollvm} --- компилятор, построенный на инфраструктуре LLVM.

На первый взгляд, \texttt{gollvm} мог бы полностью снять необходимость разработки собственного векторизатора для Go:
LLVM содержит зрелые и активно развиваемые проходы цикловой и SLP-векторизации,
способные генерировать высокоэффективный SIMD-код.
Тем не менее использование \texttt{gollvm} в качестве основного инструмента для внедрения автоматической векторизации
в экосистему Go сопряжено с рядом проблем, делающих его неприемлемым на практике:

\begin{itemize}
    \item \textbf{Несовместимость с рантаймом Go.} Стандартный рантайм Go тесно завязан на конкретные свойства кода: точный сборщик мусора опирается на карты указателей, планировщик горутин ожидает определённой структуры стековых фреймов и корректного сохранения/восстановления регистров в точках синхронизации. Поэтому \texttt{gollvm} пользуется собственной версией рантайма, которая сильно уступает основной.
    
    \item \textbf{Отставание от основного компилятора и нестабильность.} \texttt{gollvm} является экспериментальным проектом
    и не входит в официальный набор инструментов Go. 
    Его развитие отстаёт от осовоной имплементации: внедрение новых возможностей языка,
    изменений в спецификации или внутреннего устройства рантайма требует длительной адаптации.
    
    \item \textbf{Несоответствие требованию быстрой компиляции.} LLVM предоставляет богатый набор глубоких анализов и 
    агрессивных оптимизаций, но ценой значительного увеличения времени компиляции. 
    Сборка \texttt{gollvm}'ом может занимать в разы или даже на порядок больше времени по сравнению со стандартным компилятором. 
    Учитывая, что быстрота компиляции является одной из ключевых ценностей языка,
    переход на LLVM-бэкенд для массового использования маловероятен --- это противоречит ожиданиям разработчиков.
\end{itemize}

\textbf{Вывод для настоящей работы.} Таким образом, хотя \texttt{gollvm} демонстрирует потенциальные выгоды от векторизации Go-кода,
он не может служить практическим решением для большинства пользователей языка.
Единственным путём внедрения автоматической векторизации в Go является разработка соответствующих оптимизационных проходов
непосредственно внутри стандартного компилятора с полным учётом его архитектуры, ограничений и требований рантайма.
Именно это и составляет предмет данной работы. 
Далее в главе рассматривается структура и особенности стандартного компилятора Go, значимые для реализации SLP-векторизатора.

\subsection{SSA-представление}
\label{sec:go_ssa}

SSA-представление компилятора Go (пакет \texttt{cmd/compile/internal/ssa}) реализовано в виде набора связанных структур данных, ключевыми из которых являются \texttt{Block} (базовый блок) и \texttt{Value} (инструкция-значение). Каждая функция после построения SSA состоит из графа базовых блоков, внутри каждого из которых располагается линейный список значений. Значения нумеруются в порядке следования инструкций, что позволяет быстро сравнивать их позиции.

Основные свойства SSA-формы в Go:

\begin{itemize}
    \item \textbf{Явные зависимости по данным} --- каждое значение ссылается на свои аргументы. 
    Это упрощает анализ потока данных.
    \item \textbf{Единый тип значения} --- каждое значение имеет конкретный тип: указатели, целые числа, 
    числа с плавающей точкой и составные типы.
    \item \textbf{Моделирование памяти} --- операции, взаимодействующие с памятью (загрузки, сохранения, вызовы функций), 
    принимают и возвращают специальное значение типа \texttt{TypeMem}, 
    образуя цепочку состояний памяти. Это гарантирует корректный порядок операций.
    \item \textbf{Коды операций} --- множество операций разделено на архитектурно-независимые и архитектурно-зависимые. 
    В процессе компиляции обобщённые операции понижаются до машинно-зависимых.
    \item \textbf{Отсутствие данных для циклового анализа} --- SSA-форма не содержит информации о циклах, 
    индуктивных переменных или регионах; вся информация о структуре программы представлена графом базовых блоков.
\end{itemize}

Эти свойства делают SSA-представление удобным для реализации SLP-векторизации: 
зависимости данных прозрачны, типы известны статически, а память контролируется явными аргументами, позволяющими консервативно,
но уверенно гарантировать отсутствие конфликтов.

\subsection{Оптимизационный конвейер}
\label{sec:go_opt_pipeline}

Оптимизационный конвейер компилятора Go состоит из фиксированной последовательности проходов,
выполняемых над SSA-представлением функции. 
К числу основных проходов относятся (в порядке их исполнения):

\begin{itemize}
    \item \texttt{deadcode} --- устранение мёртвого кода (удаление инструкций, результаты которых не используются).
    \item \texttt{prove} --- устранение лишних проверок, в том числе проверок выходов за границы массивов.
    \item \texttt{memcombine} --- объединение соседних узких загрузок/сохранений в широкие в тех случаях, когда адреса смежны и выровнены.
    \item \texttt{lower} --- понижение до машинных операций (замена обобщённых \texttt{Op}-кодов на архитектурно-специфичные).
    \item \texttt{regalloc} --- распределение регистров (замена виртуальных регистров физическими).
\end{itemize}

Проходы выполняются в строго определённом порядке, и каждый из них может изменять SSA-граф. 
Для внедрения SLP-векторизации важно выбрать этап, на котором отсутствуют машинно-зависимые операции, память явно представлена, 
а избыточные проверки уже удалены. В работе выбран этап до \texttt{lower}.

\subsection{Особенности SSA, значимые для SLP-векторизации}
\label{sec:ssa_slp_features}

SSA-форма имеет ряд особенностей, которые определяют архитектуру будущего прохода:

\begin{itemize}
    \item \textbf{Явное кодирование зависимостей по памяти.} Каждая операция с памятью явно принимает и возвращает 
    \texttt{TypeMem}, что позволяет консервативно проверять независимость инструкций: две нагрузки могут 
    быть упакованы только в том случае, если они разделяют общее состояние памяти (или связаны допустимой цепочкой зависимостей).
    \item \textbf{Наличие информации о типах.} Для каждого значения известен его тип и размер.
    Это позволяет быстро фильтровать кандидатов по совместимости (одинаковая разрядность) и вычислять ширину будущего векторного значения.
    \item \textbf{Отсутствие сложных косвенных указателей.} Благодаря тому, что Go не поддерживает арифметику
    указателей общего назначения, анализ смежности адресов сводится к сравнению базовых указателей и константных смещений,
    что существенно проще общего случая C/C++.
    \item \textbf{Ограниченный набор подлежащих векторизации операций.} Только базовые арифметические,
    логические операции и обращения с памятью.
    \item \textbf{Активное развитие поддержки SIMD-инструкций.} На момент выполнения работы компилятор Go
    уже содержал встроенную поддержку векторных инструкций для нескольких архитектур.
    В частности, для x86-64 реализованы операции, использующие расширения SSE и AVX.
    Для ARM64 поддерживаются инструкции NEON.
    В стадии активной разработки находятся поддержка масштабируемых векторных расширений ARM SVE и векторных инструкций WebAssembly.
\end{itemize}

Эти свойства требуют аккуратного учёта при проектировании.

\subsection{Архитектурные ограничения для внедрения SLP}
\label{sec:go_limitations2}

Существующая архитектура компилятора Go накладывает ряд ограничений.

\subsubsection{Явное моделирование памяти}

Существенной особенностью SSA-представления в компиляторе Go является явное представление состояния памяти через специальный тип \texttt{TypeMem}. 
В отличие от классических SSA-представлений, где операции доступа к памяти часто рассматриваются как побочные эффекты, 
в SSA-представлении компилятора Go память представлена явно в графе зависимостей данных.

Каждая операция чтения из памяти принимает текущее состояние памяти в качестве одного из аргументов. 
Операции записи формируют новое состояние памяти, которое затем используется последующими операциями. 

Например:

\begin{verbatim}
m0 = InitMem              // создание начального состояния памяти
v0 = ConstInt 0
v1 = Load ptr1, m0
m1 = Store ptr2, v0, m0   // изменение состояния памяти
v2 = Load ptr3, m1
\end{verbatim}

Для реализации SLP-векторизации данная особенность имеет двоякое значение.

С одной стороны, наличие явных зависимостей по памяти упрощает проверку корректности преобразований, 
поскольку позволяет непосредственно анализировать порядок операций доступа к памяти.

С другой стороны, объединение нескольких операций чтения или записи в одну векторную операцию требует строгого соблюдения семантики цепочки состояний памяти.
Например, две операции загрузки могут быть объединены только в случае отсутствия промежуточных модификаций памяти, способных изменить наблюдаемое значение. 

Таким образом, работа с типом \texttt{TypeMem} является важной особенностью для реализации SLP-векторизации в компиляторе Go 
и требует специального учёта при анализе допустимости преобразований.

\subsubsection{Ограниченная аналитическая инфраструктура}

Компилятор Go не содержит многих анализов, на которые опирается векторизатор LLVM:
\begin{itemize}
    \item отсутствует \texttt{ScalarEvolution} --- нет точного анализа выражений адресации и сравнения смещений в общем виде;
    \item отсутствует \texttt{AliasAnalysis} --- невозможно определить, пересекаются ли области памяти двух произвольных указателей;
    \item нет \texttt{LoopInfo} --- нельзя получить информацию о границах циклов и инвариантах.
    \item отсутствует модель стоимости векторных операций (TargetInfo) --- невозможно оценить оценить, когда выгоднее использовать скалярный код, 
    особенно для вычислений, требующих перекомпановки данных. 
\end{itemize}
Единственным доступным средством анализа смежности адресов является логика, реализованная в проходе \texttt{memcombine}. 
Она опирается на разбор указателей вида «база + индекс + смещение».
SLP-векторизатор вынужден либо повторно использовать эту логику, либо реализовывать собственные упрощённые проверки.

\subsubsection{Требования к времени компиляции}

Одним из ключевых принципов разработки компилятора Go является сохранение высокой скорости сборки. Любой новый проход, в том числе векторизатор, обязан вносить минимальные накладные расходы.
Это означает, что алгоритмы поиска и расширения групп должны быть ограничены по сложности (линейной или почти линейной относительно размера блока), 
а дорогостоящие рекурсивные обходы значительно ограничены.
Таким образом, проектируемый векторизатор должен находить баланс между глубиной анализа и временем компиляции, сознательно отказываясь от ряда возможностей ради сохранения быстродействия.

\subsubsection{Особенности рантайма языка}
\label{sec:go_runtime}

Рантайм языка Go накладывает дополнительные ограничения, которые необходимо учитывать при внедрении SIMD-операций в общий компиляторный пайплайн.
В первую очередь это касается взаимодействия с сборщиком мусора и механизмом переключения горутин.

\textbf{Сборщик мусора и указатели в SIMD-регистрах.} Go использует точный сборщик мусора, который для своей работы требует в
любой точке синхронизации (safe-point) полной информации о живых указателях на кучевые объекты. 
Традиционно такие указатели хранятся в регистрах общего назначения или на стеке, где компилятор может сгенерировать для них специальные карты. 
SIMD-регистры, как правило, не предназначены для хранения скалярных указателей и не поддерживают генерацию соответствующих метаданных в текущей инфраструктуре компилятора Go. 
Как следствие, векторизованный код не должен содержать указатели на Go-объекты внутри SIMD-регистров в моменты, где возможна остановка для сборки мусора. 
На практике это означает, что SLP-векторизатору безопасно работать только с типами, не содержащими указателей.

\textbf{Планировщик горутин и сохранение SIMD-контекста.} Рантайм Go осуществляет кооперативное и некооперативное переключения горутин, 
в ходе которых должен быть сохранён и восстановлен полный контекст исполнения, включая значения всех регистров.
Современные версии компилятора Go уже включают в сохраняемый фрейм горутины состояние SSE/AVX/NEON регистров (например, \texttt{XMM}, \texttt{YMM} на x86-64, \texttt{V}-регистры на ARM64), 
что добавляет небольшие, но ненулевые накладные расходы при каждом переключении. Это необходимо принимать во внимание при оценке общей эффективности векторизации. 
Однако для большинства вычислительно насыщенных участков, где векторизация приносит наибольший эффект, переключения горутин редки, и данными накладными расходами можно пренебречь.

Таким образом, особенности рантайма Go задают некоторые ограничения применимости SLP-векторизатора: обработка данных без указателей, 
ограничение векторизации между местами потенциального кооперативного вытеснения (например, вызовами функций).

\subsection{Выбор точки интеграции}
\label{sec:integration_point}

В рамках данной работы проход SLP-векторизации размещается непосредственно перед этапом \texttt{lower}, преобразующим машинно-независимую SSA-форму в машинно-зависимую.

Выбор данной точки интеграции обусловлен особенностями организации оптимизационного конвейера компилятора Go.
К моменту выполнения прохода уже завершены основные SSA-оптимизации, включая локальные упрощения и устранение избыточных вычислений.
Это позволяет анализировать относительно стабилизированную структуру вычислительного графа. Одновременно с этим программа всё ещё представлена в
достаточно высокоуровневой SSA-форме, сохраняющей семантическую структуру операций и типов.

Размещение прохода до понижения SSA-формы к машинно-зависимой позволяет реализовать общий векторизатор, используя уже существующую инфраструктуру
понижения обобщённых SIMD-операций в машинные инструкции целевой архитектуры.

\subsection{Выводы по главе}
\label{sec:go_conclusion}

В данной главе была рассмотрена архитектура SSA-бэкенда компилятора языка Go и проведён анализ её особенностей с точки зрения реализации автоматической SLP-векторизации.

Показано, что используемое в компиляторе SSA-представление предоставляет удобную основу для локального анализа зависимостей и выявления векторизуемых
ациклических участков кода. Этому способствует явное представление вычислительных зависимостей между операциями, строгая типизация SSA-инструкций, 
а также использование специального типа \texttt{TypeMem}, обеспечивающего явное моделирование состояния памяти в графе вычислений.

Установлено, что наличие развиваемой инфраструктуры SIMD-операций на уровне SSA позволяет реализовывать автоматическую
векторизацию без необходимости создания отдельного промежуточного представления для векторных вычислений.

Вместе с тем выявлен ряд архитектурных ограничений, влияющих на проектирование алгоритма.
К ним относятся ограниченность существующей аналитической инфраструктуры и необходимость сохранения малого времени компиляции.

На основании проведённого анализа в качестве точки интеграции SLP-векторизующего прохода выбран этап, непосредственно предшествующий этапу понижения 
до машинно-зависимой SSA-формы. Такое размещение позволяет выполнять анализ над высокоуровневым SSA-представлением и одновременно использовать существующий
механизм последующего преобразования SIMD-операций в машинные инструкции целевой архитектуры.

Полученные результаты определяют основные требования к проектируемому алгоритму SLP-векторизации, рассматриваемому в следующей главе.